\chapter{Silicon Host and Dopant Research Questions}
\label{ch:physical-problem}

\textbf{Status: expository synthesis of the accepted physical specification;
research outcomes remain proposed work.}

\section{The parent description}

The program asks when a first-principles Kohn--Sham description of silicon and
substitutional dopants can be replaced by a smaller model without losing the
operator information needed for its declared use. The question is not whether
one model is universally ``correct.'' It is whether a stated reduced-model
class preserves stated features of a stated parent problem within separately
declared errors.

The accepted physical conventions are owned by
\path{specification/ksdft2Effmass.physical-specification.v1.md}. This chapter
organizes those conventions as a research narrative and introduces no new
physical branch.

The parent is a self-consistent Kohn--Sham density-functional calculation in
the Born--Oppenheimer fixed-ion framework~\cite{hohenberg1964,kohn1965}. Its
output is a one-particle Kohn--Sham operator and associated finite numerical
representations. It is not the complete interacting many-body Hamiltonian, and
its eigenvalues are not identified here with the complete many-body excitation
spectrum.

For the version-one program, exchange and correlation are described with the
PBE generalized-gradient approximation~\cite{perdew1996}, and the
electron--ion interaction is represented by compatible optimized
norm-conserving Vanderbilt pseudopotentials from the approved PseudoDojo
family~\cite{hamann2013,vansetten2018}. These choices define the parent against
which later reductions are measured. In particular, disagreement between this
parent and experiment is a parent-model question, not automatically an error
introduced by tight-binding or effective-mass reduction.

\section{Pristine silicon host}

The host is crystalline silicon in the diamond structure: an fcc Bravais
lattice with a two-atom basis. Primitive-cell and conventional-cell
coordinates may both be useful, but every retained representation must
identify which lattice vectors, axes, basis sites, and reciprocal-coordinate
convention it uses.

The immediate bulk pilot is scalar relativistic, non-spin-polarized, and does
not include spin--orbit coupling. Its production lattice constant is to be
obtained from a zero-pressure PBE bulk relaxation using the selected silicon
pseudopotential. Until that calculation and its convergence record exist, no
numerical lattice constant is a calculated result of this project.

The pilot targets a band-edge subspace sufficient for an orthogonal
$sp^3s^{\ast}$ reduction, including the valence and low-conduction states
needed to describe the indirect Kohn--Sham gap and electron effective masses.
The exact band count, projection choices, and Wannier windows belong to later
methodological records. They are not fixed by the phrase ``band-edge
subspace'' alone.

The declared bulk observables are the indirect Kohn--Sham gap, the
conduction-valley position, longitudinal and transverse electron effective
masses, and withheld-point band errors. These quantities serve distinct roles:
they characterize the selected parent, test numerical convergence, and later
evaluate reduced representations. Passing one role does not imply passing the
others.

\section{Compatible host and dopant branches}

Each doped calculation uses a periodic supercell built from the frozen host
lattice vectors, with one substitutional dopant and a matched pristine
comparison supercell. In the primary dopant branches, internal atomic
coordinates are relaxed while the approved bulk lattice vectors remain fixed;
unrelaxed structures are retained only as diagnostics. Geometry and site
correspondence must be identified before pristine and doped operators are
compared.

The immediate non-spin-polarized, non-SOC bulk pilot does not by itself supply
every host reference required by the dopant program. A phosphorus comparison
requires a compatible bulk reference with documented spin and SOC conventions.
Final boron acceptor and continuum claims require a matching fully relativistic,
noncollinear SOC bulk reference. Operators from these different spin spaces are
not directly subtractable without an explicit identification.

\section{Substitutional phosphorus}

The phosphorus branch replaces one silicon atom at a diamond-lattice site by
phosphorus in a periodic supercell. Its primary version-one charge state is
neutral $\mathrm{P}_{\mathrm{Si}}^{0}$. With the accepted pseudopotential
valences, this system has one additional valence electron relative to the
matched pristine silicon supercell and therefore an odd valence-electron count.
The primary branch is consequently scalar relativistic, collinearly spin
polarized, and non-SOC.

The neutral supercell contains both the donor-core perturbation relative to
silicon and the additional electron that occupies and screens that
perturbation. It must not be relabelled silently as an isolated ionized donor.
A charged $\mathrm{P}_{\mathrm{Si}}^{+}$ calculation, or an extraction intended
to isolate an ionized-core potential, requires a separately specified
controlled branch.

The eventual phosphorus questions concern donor energies relative to a
compatible conduction-band reference, valley composition, localization or
envelope diagnostics, and the spatial structure of an aligned impurity
operator. Each finite supercell represents a periodic dopant array. An
isolated-donor interpretation requires supercell-size evidence rather than the
presence of a single substituted atom in the simulation cell.

\section{Substitutional boron}

The boron branch replaces one silicon atom by boron and uses neutral
$\mathrm{B}_{\mathrm{Si}}^{0}$ as its primary version-one charge state. The
accepted valence convention gives one fewer valence electron than in the
matched pristine supercell and again produces an odd electron count.

For final acceptor and continuum claims, the physical specification requires a
fully relativistic, noncollinear spinor calculation with spin--orbit coupling.
A scalar-relativistic non-SOC calculation may be retained only as an explicitly
labelled method-development branch. This distinction prevents an early
simplification from being presented as the final physical treatment of the
spin--orbit-coupled valence manifold.

The eventual boron questions concern acceptor energies relative to a
compatible valence-band reference, valence-subspace fidelity, localization or
envelope diagnostics, and the radial structure of the aligned impurity
operator. Charged $\mathrm{B}_{\mathrm{Si}}^{-}$ calculations are deferred to
a separate controlled branch.

\section{From parent models to reduced models}

The program studies a sequence of representations rather than one direct leap
from density-functional theory to continuum theory:
\begin{enumerate}
    \item a converged Kohn--Sham parent problem;
    \item a retained projected or Wannier subspace;
    \item an aligned finite operator representation;
    \item a constrained lattice model, initially an orthogonal
          $sp^3s^{\ast}$ hierarchy for bulk silicon;
    \item an extracted and reduced impurity operator; and
    \item where justified, a continuum effective-mass description.
\end{enumerate}

Projection, localization, alignment, parameter fitting, truncation, and
continuum approximation are different operations. Their errors and assumptions
must therefore remain identifiable. Wannier localization alone is not a
low-rank approximation, and a good spectral fit alone does not establish
agreement of aligned operator matrix elements.

\section{Research questions}

Within these boundaries, the long-form study is organized by the following
questions:
\begin{enumerate}
    \item Which finite subspace and representation preserve the declared bulk
          Kohn--Sham features with controlled numerical error?
    \item Can a prescribed tight-binding model class satisfy both spectral and
          aligned-operator criteria, and what is the smallest tested class that
          does so?
    \item What basis, gauge, geometry, unit, and energy-reference
          identifications are required before pristine and doped operators can
          be compared?
    \item After those identifications, what portion of the
          substitutional-dopant perturbation can be represented by nested local
          or nonlocal lattice-model classes?
    \item Under separately defined acceptance rules, when does a continuum
          effective-mass description preserve the declared donor or acceptor
          observables?
\end{enumerate}

These are proposed research questions, not statements of completed findings.
Their answers require retained calculations and evidence produced by the
applicable downstream tasks.

\section{Interpretive limits}

A finite periodic supercell is not an isolated impurity. A Kohn--Sham band gap
is not an exact quasiparticle gap. Spectral agreement is not operator equality.
Matrix subtraction is not physically meaningful before common-space and
energy alignment. Software execution and passing tests establish neither
numerical convergence nor scientific validation unless the corresponding
evidence and acceptance rules are present.

These limits are part of the scientific problem rather than editorial caveats:
they determine which comparisons are well posed and which conclusions the
resulting hierarchy can support.
